\section{The compact base of the global theta comparison}
\label{sec:sr-compact-global-comparison}

Fix an actual triple from the preceding local construction, put
\(d=a+a'\), \(\nu=\alpha/\beta\), and suppose \(p=q=0\).
Thus \(V\), \(V'\), and \(iW\) are definite at every infinite
place, of signatures \((a,0)\), \((0,a')\), and \((d,0)\).
The inequalities in the triple imply \(a'\geq2\), \(a>a'\),
\(a\equiv a'\pmod2\), and
\[
       \kappa_w=m_{\alpha,w}-m_{\beta,w}>d
                  \quad\hbox{at every infinite place }w.
\]
All measures below are the measures of the actual models.
In particular, replacing an invariant measure by probability
measure is always accompanied by its displayed volume factor.
Write
\[
 [H]=H(F)\backslash H(\mathbb A_F),\qquad
 I_\Phi(h)=\int_{[U(V)]}\theta_\Phi(g,h)\,dg,\qquad
 I'_{\Phi'}(h)=\int_{[U(V')]}\theta'_{\Phi'}(g,h)\,dg.
\]
These three adelic quotients are compact: a rational isotropic
vector would remain isotropic at a definite infinite place,
so each group is anisotropic over \(F\); reduction theory for
anisotropic reductive groups gives compactness. Consequently
ordinary theta convergence, differentiation on compact sets,
and Fubini apply to the Schwartz vectors used here.

\subsection{Factoring the actual theta periods through the common module}

We first supply the factorization needed for the global argument;
the existence of common local quotient maps alone does not imply it.

\begin{srlemma}
Let \(Q\) be a finite-length admissible smooth representation of
a totally disconnected group. Suppose that its maximal semisimple
quotient is one irreducible representation \(\pi\), with
multiplicity one. If a map from \(Q\) to a unitary representation
has an invariant positive semidefinite pullback form, its kernel
contains the kernel of \(Q\twoheadrightarrow\pi\).
\end{srlemma}

\begin{proof}
Pass to the image, with its positive definite form. Every invariant
submodule of this image has an invariant orthogonal complement
inside the algebraic image. Indeed, for a vector fixed by a compact
open subgroup \(K\), its orthogonal projection to the Hilbert closure
of a submodule belongs to the closure's \(K\)-fixed space.
That space equals the submodule's \(K\)-fixed space: averaging
approximants over \(K\) proves density, and admissibility makes the
latter space finite dimensional and hence closed. Thus the
projection still belongs to the algebraic submodule. Orthogonal
complements are invariant because the action is unitary.
Induction on length shows that the image is semisimple.
Every semisimple quotient of \(Q\) factors through its maximal
semisimple quotient, proving the assertion, including the
possibility of zero image.
\end{proof}

At a finite place take the quotient of the local Weil module
by the first unitary group's augmentation space. It is the big
theta lift of the trivial representation for the \emph{prescribed}
splitting. It has finite length, and its maximal semisimple quotient
is irreducible. For a nonsplit place these are the finite-length
statement and Theorem~1.2 of Gan--Takeda,
\emph{A proof of the Howe duality conjecture},
\url{https://arxiv.org/pdf/1407.1995}; the ground field has
characteristic zero, and the pair is unitary, so their hypotheses
apply. At a split place use the type-II result of
M\'inguez, \emph{Correspondance de Howe explicite: paires duales
de type II}, Corollary~6.3 and the big-lift finite-length statement
in its introduction,
\url{https://arxiv.org/pdf/0709.4384}.
Its model
\[
  (g,h)f(Y)=|\det g|^{-d/2}|\det h|^{n/2}f(g^{-1}Yh)
\]
is our split model after \(g\mapsto(g^{\mathsf t})^{-1}\)
and the two prescribed determinant characters; these automorphisms
and character twists preserve both finite length and the
multiplicity-one assertion. The real and finite model
identifications were established before constructing the maps,
so these statements concern the actual Weil actions.

For a fixed vector at all other places, the map
\(\phi_v\mapsto I_{\phi_v\otimes\Phi^v}\) takes values in
\(L^2([U(W)])\). First-group averaging kills the augmentation
space, and second-group translation preserves its \(L^2\) form.
The lemma therefore proves that this map kills the kernel of
the actual irreducible local quotient map. The same argument
applies to the primed model.

At infinity all the groups in this base case are compact.
Averaging the first group on polynomial Fock vectors is an
algebraic projection onto its invariant part: each such vector
has a finite-dimensional compact-group orbit. The Cauchy
decomposition of polynomials on a matrix space identifies each
prescribed first determinant type with one second Schur module.
The determinant degrees are the rectangular degrees computed
in the real matching construction. Hence the first-invariant
part is exactly the common irreducible second module there,
and averaging kills precisely the kernel of its quotient.
This proves the same factorization at infinity without any
claim about arbitrary extensions of the real maps off the
polynomial vectors.

Let
\[
             \Pi=\Pi_\infty\otimes
                       \bigotimes_{v<\infty}'\Pi_v
\]
use the common finite quotients and their common distinguished
vectors outside \(T\). Exactness of tensor products of complex
vector spaces now gives maps
\[
        j,j':\Pi\longrightarrow C([U(W)]),\qquad
        I_\Phi=j(r\Phi),\quad I'_{\Phi'}=j'(r'\Phi').
\]
To see the kernel assertion explicitly, write a vector as a
finite sum of pure tensors and enlarge the finite exceptional
set until every tensor is distinguished outside it. The kernel
of the quotient map on that finite tensor product is the sum of
the kernels in its individual factors. Each summand has just
been killed by the theta period. The transition maps preserve
the common distinguished vectors, so the factorization does
not depend on the exceptional set.

The module \(\Pi\) is irreducible and admissible. For irreducibility,
a nonzero vector involves finitely many factors; the density
theorem for finitely many irreducible modules isolates a pure
tensor, after which the individual group actions generate every
finite tensor product. At almost all places the distinguished
vector generates the irreducible local module. Admissibility
follows from local admissibility and the one-dimensional
hyperspecial fixed lines outside a finite set. The module has
countable complex dimension, so its endomorphisms are scalars
by the countable-dimensional Schur argument proved in the real
matching construction.

We shall use the following slightly stronger consequence.
If \(j'\ne0\), every invariant sesquilinear form on \(\Pi\)
is a scalar multiple of
\(\gamma(u,v)=\langle j'(u),j'(v)\rangle\).
Indeed, \(j'\) is injective. Given another such form \(B\) and
a \(K\)-fixed vector \(u\), invariance gives
\(B(u,v)=B(u,e_Kv)\). The finite-dimensional positive space
\(\Pi^K\) has a unique vector \(Tu\) representing this functional
with respect to \(\gamma\). This represents it on all of \(\Pi\)
by averaging; hence it is independent of the chosen \(K\).
The resulting linear map \(T\) commutes with the group and is
a scalar. Here \(K\) is compact open in the finite adelic group;
\(\Pi_\infty\) is finite dimensional because \(U(W)_\infty\)
is compact in this base case, so \(\Pi^K\) is still finite
dimensional. This argument also applies to the mixed forms
defined by \(j\) and \(j'\).

\subsection{The local matrix-functional calculation}

The following elementary distribution calculation will turn a
theta back-lift into a constant. It applies to algebraic linear
functionals on locally constant compactly supported functions;
no continuity of such a functional is assumed.

\begin{srlemma}
Let \(k\) be a nonarchimedean local field with uniformizer \(\pi\).
Let \(\xi:k^\times\to\mathbb C^\times\) be smooth and suppose
\(\xi(\pi)\notin\mathbb R_{>0}\).
If
\[
 \ell(f(\,\cdot\,h))=\xi(\det h)\ell(f)
       \quad(h\in\operatorname{GL}_d(k))
\]
on \(\mathcal S(\operatorname{Mat}_{n,d}(k))\), then
\(\ell=0\) when \(n<d\). If \(n=d\), then also
\[
                \ell(f(g\,\cdot\,))=\xi(\det g)\ell(f).
\]
\end{srlemma}

\begin{proof}
Filter the matrix space by rank. Within rank \(r\), right
translation is transitive on the matrices with a specified
column space. After choosing a basis of this column space,
the stabilizer of \([1_r\ 0]\) is
\[
       H_r=\left\{\begin{pmatrix}1_r&0\\ C&D\end{pmatrix}
                     :D\in\operatorname{GL}_{d-r}(k)\right\}.
\]
The elementary integration criterion for an equivariant
functional on \(H_r\backslash\operatorname{GL}_d(k)\)
requires the restriction of \(\xi\circ\det\) to \(H_r\)
to equal the quotient of the two modular characters.
This criterion follows directly by writing a test function
as an integral over a compact open subgroup and then changing
variables by an element of \(H_r\); conversely the same formula
defines the functional when the characters agree. Thus it is
valid for the algebraic dual of \(\mathcal S\) as well.
For \(r<d\), use \(D=\operatorname{diag}(\pi,1,\ldots,1)\).
The modular-character quotient is a positive real power of
\(|\pi|\), whereas \(\xi(\pi)\) is not positive real.
There is consequently no functional on this orbit.
To pass from the individual column spaces to the whole rank
stratum, use precisely the following localization principle:
if \(q:Z\to T_0\) is a continuous map of locally compact,
totally disconnected Hausdorff spaces, a locally compact
totally disconnected group acts continuously on \(Z\)
preserving every fiber, and \(\mu\) is a continuous character,
then vanishing of all \(\mu\)-equivariant distributions on
every fiber implies vanishing on \(Z\).
This is Bernstein localization, in the form of Corollary~2.1
of Aizenbud--Gourevitch--Rallis--Schiffmann,
\emph{Multiplicity one theorems}, Annals of Mathematics
172 (2010), \S2:
\url{https://annals.math.princeton.edu/wp-content/uploads/annals-v172-n2-p14-p.pdf}.
Apply it to the rank-\(r\) stratum and the map sending a
matrix to its column space in
\(\operatorname{Gr}_r(k^n)\). Both spaces have all the stated
topological properties and are countable at infinity.
The column-space map is continuous: on a chart with an
invertible \(r\)-minor its graph coordinates are obtained
by multiplying the remaining rows by that minor's inverse.
Right \(\operatorname{GL}_d(k)\) preserves the map, and
each fiber is exactly the homogeneous orbit calculated
above. The character \(\xi\circ\det\) is smooth and hence
continuous. Thus every deficient-rank stratum has zero
equivariant distributions.

There is no unmentioned continuity assumption on \(\ell\).
For any of these second-countable totally disconnected
spaces, the Bruhat space \(\mathcal S\) is the inductive
limit of its finite-dimensional spaces with prescribed
compact support and prescribed compact-open subdivision.
Every algebraic linear functional is continuous in this
topology, since its restrictions to all these finite-dimensional
spaces are continuous. It is therefore exactly a distribution
to which localization applies.
Finally the rank-at-most-\(r\) subsets are closed.
Restriction of test functions to a closed subset is onto:
a locally constant function of compact support there extends
on finitely many compact-open neighborhoods and can be
made zero off their union. A function vanishing on that
subset has compact support away from it after such a
subdivision. This gives the exact open/closed test-function
sequence. Starting with the largest rank and descending
therefore kills all strata with \(r<d\).

If \(n<d\), these are all strata. If \(n=d\), only the open
orbit \(\operatorname{GL}_d(k)\) remains. On that orbit an
equivariant functional is a scalar multiple of
\[
                  f\longmapsto
                    \int_{\operatorname{GL}_d(k)}
                         f(Y)\xi(\det Y)^{-1}\,d^\times Y .
\]
On compact support inside the open orbit this is an ordinary
finite integral and has the asserted left transformation law.
Subtracting that left law from \(\ell\) gives a functional
supported on deficient rank, already shown to vanish. No
extension of the displayed integral across singular matrices
is used.
\end{proof}

Suppose now \(u\in\Pi\) and put
\(F_u=sj(u)+tj'(u)\), for fixed scalars \(s,t\). For a Schwartz
vector \(\Psi\) in either model form the back-lift
\[
     T_{\Psi,V}F_u(g)=
         \int_{[U(W)]}F_u(h)\,
                          \overline{\theta_\Psi(g,h)}\,dh ,
\]
and define \(T_{\Psi',V'}\) similarly. At a good split finite
place both models are the actual matrix models. To obtain
the required local functional even when \(s\ne0\), express
\(u\) as a finite sum of tensors \(u_v\otimes u^v\) and treat
each summand. The complementary local quotient
\[
 r'_v:\mathcal S(\operatorname{Mat}_{a',d}(F_v))
                                      \twoheadrightarrow\Pi_v
\]
is surjective. For fixed \(u^v\), define the \emph{linear}
automorphic-valued map
\[
 C_v(\phi')=
    s\,j(r'_v(\phi')\otimes u^v)
           +t\,j'(r'_v(\phi')\otimes u^v).
\]
It is equivariant for the actual second local Weil action
because both \(j\) and \(j'\) are maps of the common module.
Pair \(C_v(\phi')\) with the unprimed local theta test
\(\psi_v\), freezing every other component of that theta
test. After conjugating \(\psi_v\) this is a bilinear
functional on one \(a\)-row and one \(a'\)-row matrix
space. Its total row count is therefore exactly \(d\),
irrespective of \(s,t\); it is not a claim restricted to
a product of one primed and one unprimed theta kernel.
For the opposite back-lift use instead the onto quotient
\(r_v:\mathcal S(\operatorname{Mat}_{a,d}(F_v))
\twoheadrightarrow\Pi_v\), and pair with the \(a'\)-row
primed theta test. The same argument applies with the
two row blocks interchanged.
Simultaneous right translation of both matrices acts
on the first of these bilinear functionals by
\[
                 \xi(\det h)=
                    \nu_v(\det h)^{-1}|\det h|^{-d/2}.
\]
Here the pairing is bilinear after complex conjugating its
unprimed test vector. The source of this character can be
checked without any convention about induction: the two
second-group operators are
\[
 \beta_v(\det h)|\det h|^{a/2}f(Yh),\qquad
 \alpha_v(\det h)|\det h|^{a'/2}f'(Y'h).
\]
Invariance of the adelic integral divides by
\(\alpha_v\overline{\beta_v}|\det h|^{d/2}\), giving exactly
the displayed \(\xi\).

The good split non-torsion calculation gives
\(\nu_v(\pi)\ne1\). Since it is unitary,
\(\xi(\pi)\notin\mathbb R_{>0}\).
The square-matrix conclusion of the lemma therefore moves
left translation in the first \(a\) rows through the pairing.
The actual unprimed first operator is
\[
        \nu_v(\det g)^{-1}|\det g|^{d/2}f(g^{\mathsf t}Y).
\]
Complex conjugation of this scalar cancels the left
transformation factor
\(\nu_v(\det g)^{-1}|\det g|^{-d/2}\) from the matrix lemma.
Thus \(T_{\Psi,V}F_u\) is right invariant under \(U(V)(F_v)\).
Interchanging the two models replaces \(\nu_v\) by its inverse
and proves the corresponding invariance on \(U(V')(F_v)\).
Finite sums of local pure tensors give the statement for all
the polynomial adelic vectors. In particular this holds at
every place outside \(T\) lying over a rational prime which
splits completely in \(E\).

\subsection{From good-place invariance to a constant}

For completeness the density assertion needed here is stronger
than strong approximation for the derived group alone.

\begin{srlemma}
Let \(E/F\) be an actual CM extension, with \(E/\mathbb Q\)
Galois, and let \(H=U(X)\), \(\dim_E X\geq2\).
A continuous left \(H(F)\)-invariant function which is right
invariant under \(H(F_v)\) at every \(v\) outside a finite set
and above a rational prime splitting completely in \(E\)
is constant.
\end{srlemma}

\begin{proof}
Choose one such split place. The simply connected group
\(SU(X)\) is noncompact there, since it is
\(\operatorname{SL}_{\dim X}(F_v)\).
Strong approximation away from that place shows that the
function is invariant under \(SU(X)(\mathbb A_F)\).
The determinant quotient is the norm-one torus. It has a
continuous adelic section preserving rational points and local
supports: choose a rational nonisotropic line and let a
norm-one scalar act on that line and as identity on its
orthogonal complement. It remains to prove density on this torus.

Its rational adelic quotient is compact. If the closed subgroup
generated by rational points and the specified local norm-one
groups were proper, a character of its compact abelian quotient
would separate a point. Pull this character back along
\(z\mapsto z/\bar z\) on \(\mathbb A_E^\times\).
Adelic Hilbert 90 makes the pullback nontrivial. It is a
unitary Hecke character \(\chi\), trivial on every local group
at all degree-one primes of \(E\), apart from finitely many:
because \(E/\mathbb Q\) is Galois, a degree-one unramified
prime lies over a completely split rational prime.

For real \(s>1\), the Euler product then gives
\[
              \log|L_E(s,\chi)|=\log\zeta_E(s)+O(1)
                                      \quad(s\downarrow1).
\]
Indeed degree-one local factors outside the exceptional set
are the zeta factors. Prime ideals of residue degree at least
two contribute a uniformly convergent sum, bounded by a
constant times \(\sum_p p^{-2s}\); the higher powers in all
logarithmic Euler factors are bounded in the same manner.
The right side tends to infinity, whereas a nontrivial unitary
Hecke character has a finite value at \(1\).
This contradiction proves density and then constancy.
The classically known analytic statement used in this argument
is ordinary Hecke continuation at \(1\); no zero-free assertion
in the critical strip is used.
\end{proof}

Applying the lemma to the preceding back-lifts proves
\[
             T_{\Psi,V}F_u(g)=T_{\Psi,V}F_u(1),\qquad
             T_{\Psi',V'}F_u(g)=T_{\Psi',V'}F_u(1).
\tag{*}\label{eq:sr-base-backlift-constant}
\]
The first ranks are at least two by the actual triple
inequalities. Continuity follows by dominated convergence
on the two compact adelic quotients. The Galois hypothesis
is a field hypothesis of the actual base.

\subsection{Nonzero theta periods from explicit doubling tests}

Here is the nonvanishing input, including the Euler factors
and the bad-place choices. Let \(X\) be either \(V\) or \(V'\),
of rank \(n\), and keep the definite \(W\) of even rank \(d>n\).
The splitting character for the first group in this pair is
\(\chi=\nu^{-1}\) or \(\nu\), respectively.
It is unitary, trivial on \(\mathbb A_F^\times\), and has
infinite absolute weights \(\kappa_w>d\).

For every integer \(j\geq-d/2\),
\[
                            L_E(j,\chi)\ne0 .
\tag{**}\label{eq:sr-base-critical-nonzero}
\]
In this formula \(L_E\) is the finite Hecke function continued
from its Euler product. To verify the assertion, its completed
function has infinite factors
\(\prod_w\Gamma_{\mathbb C}(s+\kappa_w/2)\).
For \(j\geq2\) use the absolutely convergent Euler product,
and at \(j=1\) use the classical nonvanishing theorem for a
unitary Hecke character at \(1\). For \(j\leq0\), the functional
equation moves to \(1-j\geq1\) for \(\chi^{-1}\).
Both sets of gamma factors are finite and nonzero since
\(j+\kappa_w/2>0\). The conductor power and root number are
nonzero as well. The character has nonzero infinite weight,
so its finite Hecke function is entire. This proves
\eqref{eq:sr-base-critical-nonzero}.

Choose a nonzero polynomial real vector \(f_\infty\) invariant
under the first compact group. Such a vector exists without
a global nonvanishing assumption: the onto real quotient from
local matching is first invariant, so compact averaging of
a vector with nonzero quotient has the same nonzero quotient.
For \(v\in T\), which is split in \(E/F\), use the matrix
Schr\"odinger coordinates of our model and choose
\[
 \phi_v=\mathbf1_{[\,1_n\ 0\,]
                         +\pi_v^{N_v}\operatorname{Mat}_{n,d}
                                                   (\mathcal O_v)}.
\]
Take \(N_v\) large enough that the determinant character is
trivial on
\(K_v^0=1+\pi_v^{N_v}\operatorname{Mat}_{n,n}(\mathcal O_v)\)
and that its doubled embedding is in the prescribed maximal
compact. Normalize the coordinate intertwiner by the
self-dual additive measures. Its Hilbert norm is then a
positive scalar multiple of the original model's norm.
A direct overlap calculation gives
\[
 \langle\omega_v(g,1)\phi_v,\phi_v\rangle
            =\|\phi_v\|^2\mathbf1_{K_v^0}(g).
\]
Indeed, overlap of the first \(n\) columns forces
\(g\in K_v^0\); conversely that subgroup preserves the box
and its determinant twist is trivial. A transpose or inverse
in the chosen split model gives the same subgroup.

At \(v\notin T\) use the actual self-dual vacuum.
The doubled vacuum is fixed by the prescribed doubled
maximal compact when \(\chi\) is unramified. To check the
scalar, use the Heisenberg eigenline of the tensor lattice
and its prescribed quadratic refinement. Integral Siegel
unipotents fix that line by the refinement formula; integral
Levi elements have unit determinant and hence trivial
unramified character; and the Fourier generator fixes the
self-dual lattice indicator with self-dual measure.
These elements generate the integral doubled unitary group.
At an inert place two is a unit by the base hypothesis;
at a split place the same assertion is the elementary
generation of the integral general linear group, including
residue characteristics two and three.
Transporting this calculation to the actual Darboux model
preserves the scalar. Thus its doubled matrix-coefficient
section is the spherical section with value one on the
maximal compact, multiplied by the positive squared norm
of the vacuum. All but finitely many of these norms equal
one; the exceptional positive normalization factors may
be absorbed once into the global Schwartz norm. Put
\(\Phi=f_\infty\otimes\bigotimes_v'\phi_v\).

Write \(s_0=(d-n)/2>0\), let \(\iota\) denote the actual
doubling embedding for \(X\), and let \(A\) be its positive
height, equal to one on its maximal compact. The normalized
flat section on the doubled \(X\)-group has restriction
\[
 Z\text{-integrand}(s,g)=
 A(\iota(g,1))^{s-s_0}
           \langle\omega(g,1)\Phi,\Phi\rangle.
\]
It transforms under the Siegel Levi by
\(\chi(\det t)|\det t|_E^{s+n/2}\), exactly the convention
of the classical doubling input established above.

For clarity, the unramified integral used here is
\[
 \left(\int_{U(X)(F_v)}
             f_{v,s}(\iota(g,1))\,dg\right)
 \prod_{i=0}^{n-1}
       (1-\eta_v(\pi_v)^{\,n-i-1}
                         q_v^{-(2s+i+1)})^{-1}
       =\prod_{i=0}^{n-1}
                     L_{E,v}(s+n/2-i,\chi).
\tag{***}\label{eq:sr-base-unramified-zeta}
\]
The spherical section has value one and the integral lattice
stabilizer has volume one. One may use the unramified
doubling calculation in Eischen--Harris--Li--Skinner,
\emph{p-adic \(L\)-functions for unitary groups},
\S4.2.1, equation~(53), DOI~10.1017/fmp.2020.4:
\url{https://doi.org/10.1017/fmp.2020.4}.
Here its representation is the trivial representation of
\(U(X)(F_v)\); its base-change Satake exponents are
\((n-1)/2-i\), \(0\leq i<n\), so their numerator
\(L(s+1/2,\operatorname{BC}(1)\otimes\chi)\) is exactly
our numerator. Replacing their index \(r\) by \(n-i-1\)
gives our denominator. Finally \(\chi|_{F_v^\times}=1\),
so there is no extra character in that denominator.
This is a local computation for split and inert places in
either parity; the global cuspidality and multiplicity
hypotheses appearing elsewhere in that paper are not used.

Let \(D^T_n(s)\) be the product of the displayed denominator
over \(v\notin T\). It is holomorphic and nonzero for
\(\operatorname{Re}s>0\): the smallest exponent has real
part \(2\operatorname{Re}s+1>1\).
In a sufficiently far right half-plane, absolute convergence,
Fubini on finite products and then dominated convergence give
\[
 Z(s)D^T_n(s)
      =H\prod_{i=0}^{n-1}L_E(s+n/2-i,\chi),\qquad H>0.
\tag{****}\label{eq:sr-base-global-zeta}
\]
There are no missing \(T\)-Euler factors: \(\nu\) is ramified
at every prime over each \(v\in T\), so those local Hecke
factors are one. Here \(H\) is the factor remaining \emph{after dividing
out the full unramified local integrals}. More explicitly,
choose product additive measures and product group Haar
measures, with each finite integral stabilizer of volume
one. Write the actual additive measure as \(c_S\) times
the chosen product and the actual group measure as \(c_G\)
times its product, where \(c_S,c_G>0\). With the norms
relative to the chosen additive product measures,
\[
 H=c_Sc_G\,\|f_\infty\|^2
          \operatorname{vol}(U(X)(F_\infty))
       \prod_{v\in T}\!
            \bigl(\|\phi_v\|^2\operatorname{vol}(K_v^0)\bigr)
       \prod_{v\notin T}\|\operatorname{vac}_v\|^2 .
\]
The last product has only finitely many factors different
from one. The infinite group is compact and fixes
\(f_\infty\), and the bad-place coefficient is supported
on \(K_v^0\), where its doubled height is one.
These facts account for every factor in \(H\), all positive
and independent of \(s\). At a good place the factor removed
was
\[
             J_v(s)=\int_{U(X)(F_v)}
                       f_{v,s}(\iota(g,1))\,dg ,
\]
integrated over the \emph{whole} local group, as in
\eqref{eq:sr-base-unramified-zeta}. We do not assert that
a good-place matrix coefficient has compact support.
The restricted-product Fubini identity in the convergent
half-plane gives
\(Z(s)=H\prod_{v\notin T}J_v(s)\); applying the exact
unramified calculation yields
\eqref{eq:sr-base-global-zeta}. Thus no unspecified
possibly vanishing local factor is hidden in \(H>0\).

The doubled theta identity and the anisotropic Siegel--Weil
formula from the classical-input section give a meromorphic
continuation of \(Z(s)\), regular at \(s_0\), and
\[
       I_\Phi=0\quad\Longrightarrow\quad Z(s_0)=0.
\]
Indeed the continued doubled Eisenstein convolution at this
point is a positive nonzero multiple of
\(\|I_\Phi\|^2_{L^2([U(W)])}\).
The auxiliary space is the definite skew space \(W\), of
rank \(d>n\). To apply the opposite Hermitian parity of
the same formula, multiply its form by a nonzero rational
trace-zero \(b\in E\) and the doubled \(X\)-form by \(b^{-1}\).
The two unitary groups and their tensor symplectic form
are unchanged, so the additive character, Darboux action
and self-dual measures are unchanged. This is exactly the
parity change checked in the classical-input section.
The anisotropic formula therefore has scalar \(1\) before
converting its average to the actual measure.
The doubled matrix coefficient is the one obtained from
the actual Heisenberg intertwiner, so its value at the
identity is the positive Schwartz inner product.
All changes of order here are on compact quotients.

Continue \eqref{eq:sr-base-global-zeta} to \(s_0\) in
\(\operatorname{Re}s>0\). Removing the discrete poles of
the meromorphic continuation leaves a connected domain;
the identity theorem applies, and both sides are regular
at \(s_0\). Its Hecke arguments there are the integers
\(d/2-i\), all at least \(1-d/2\); they are nonzero by
\eqref{eq:sr-base-critical-nonzero}. As \(D^T_n(s_0)\ne0\)
and \(H>0\), we obtain \(Z(s_0)\ne0\), and hence
\(I_\Phi\ne0\). This proves nonvanishing in both models.

On the primed side the good split coinvariant calculation
and the exact bar-coordinate cutoff let the nonzero vector
be replaced by one in the actual space \(S^\circ\).
The cutoff maps onto the common local quotient at \(v_1\);
its change has zero theta period by the factorization
already proved. It can be performed on finitely many pure
tensors and preserves polynomial real vectors. Thus
\(j'\ne0\) is witnessed by a vector in the space required
in the signed-data construction.

\subsection{Doubling proves equality of the two maps}

For \(u\in\Pi\) put \(F_u=sj(u)+tj'(u)\).
The invariant-form conclusion above gives a scalar
\(\lambda_{s,t}\) such that
\[
       \langle F_u,F_v\rangle
              =\lambda_{s,t}
                     \langle j'(u),j'(v)\rangle
                                     \qquad(u,v\in\Pi).
\]
Fix a primed polynomial vector \(\Psi'\).
On the doubled \(W\)-group use the box section built from
\(\Psi'\otimes\overline{\Psi'}\). For real part of the
parameter sufficiently large, ordinary Eisenstein unfolding
expresses its quadratic convolution with \(F_u\) as the
integral of that section against
\(\langle F_{\pi(h)u},F_u\rangle\).
Replace this coefficient by
\(\lambda_{s,t}\langle j'(\pi(h)u),j'(u)\rangle\).
The two meromorphic functions are therefore equal.

This unfolding has only the open orbit. In fact, a boundary
orbit would give a nonzero isotropic vector in \(W(F)\),
whereas \(iW\) is definite. The open orbit is the diagonal
copy of \(U(W)\); its stabilizer gives precisely the
convolution integral just used. Absolute convergence in
the far half-plane justifies the unfolding term by term.
The Eisenstein continuation in the classical-input section
continues the equality, with its uniform moderate-growth
bounds sufficient for integration on our compact quotient.

Evaluate at
\[
                 z_0=(a'-d)/2=-a/2 .
\]
The auxiliary space is now the definite \(V'\) of rank
\(a'<d\). The anisotropic Siegel--Weil scalar is \(2\)
and its normalized integral is the probability mean.
The actual box identity and Fubini therefore give, for
either \(F=F_u\) or \(F=j'(u)\),
\[
 \mathcal Q_{\Psi',z_0}(F)
   =\frac{2}{\operatorname{vol}[U(V')]}
             \big\|T_{\Psi',V'}F\big\|^2_{L^2([U(V')])}.
\]
The determinant character on the second doubled copy
cancels the determinant twist in the box identity.
It is identical on the two sides, as both use the same
actual primed splitting. The constant is nonzero.
Consequently
\[
 \big\|T_{\Psi',V'}F_u\big\|^2
        =\lambda_{s,t}
                    \big\|T_{\Psi',V'}j'(u)\big\|^2 .
\tag{*****}\label{eq:sr-base-backlift-norm}
\]

The mixed invariant form has a scalar \(c\) with
\[
             \langle j(u),j'(v)\rangle
                       =c\langle j'(u),j'(v)\rangle .
\]
Set \(D_u=j(u)-cj'(u)\); then \(D_u\) is orthogonal to the
entire primed image. By
\eqref{eq:sr-base-backlift-constant} its primed back-lift
is a constant function. Its average equals
\(\langle D_u,I'_{\Psi'}\rangle/
                \operatorname{vol}[U(V')]=0\);
hence the back-lift is zero. Choose \(u_0\) with
\(j'(u_0)=I'_{\Phi'_0}\ne0\) and use
\(\Psi'=\Phi'_0\). The back-lift of \(j'(u_0)\) is also
constant and its average is
\(\|j'(u_0)\|^2/\operatorname{vol}[U(V')]\ne0\).
Equation~\eqref{eq:sr-base-backlift-norm} for \((s,t)=(1,-c)\)
forces \(\lambda_{1,-c}=0\). Thus \(\|D_u\|^2=0\) for every
\(u\). Continuity on a quotient with Haar measure of full
support implies \(D_u=0\) pointwise. We have proved
\[
                    I_\Phi=c\,I'_{\Phi'}
\]
for every polynomial matched pair in the actual cutoff
space. Both nonvanishing statements show \(c\ne0\);
the primed nonzero vector shows uniqueness.

This closes the compact-base comparison, the source leaf
\(\mathrm{PolZero}\) with \(p=q=0\), including its nonzero
constant and the precise \(S^\circ\) witness. It does not
assert the full \(\mathrm{OneConst}\) for arbitrary
signatures. Propagating this base through the two
codimension-two constant-term recursions, and checking
their finite-character-independent constants, remains
the next global step. No constants in this compact-base
argument have been used as uniform family metric bounds.
